\documentclass[11pt]{article}
\usepackage{amsmath,amssymb,amsthm,booktabs}
\usepackage[margin=1in]{geometry}

\newtheorem{proposition}{Proposition}
\newcommand{\Psym}{\mathsf P_\infty}
\newcommand{\Uinf}{\mathcal U_\infty}
\newcommand{\MGamma}{\mathcal M_\Gamma}
\newcommand{\Cell}{C_\ell}

\title{Step 177: Branch B SL164.1 Attack}
\author{Six Birds Foundations III}
\date{2026-05-15}

\begin{document}
\maketitle

\section{Verdict}

\[
\boxed{\texttt{SL164\_verdict = V\_SL164\_stuck\_at\_normalization}.}
\]

The Step 173 projection identity composes cleanly with the Step 153 transport,
but the inherited records do not provide the input Mellin vector on which the
projection must act.  The deepest named gap is
\[
\boxed{
\kappa_{a,w}(\tau)
:=
\left(T_a^*K_a^\Gamma(\cdot,w)\right)(1/2+i\tau),
}
\]
equivalently the action of \(P_{\mathcal L_a^\Gamma}\) on
\(K_a^{\Gamma,\mathrm{amb}}\).

\section{T1. Inherited Records}

Step 145 defines the carrier:
\[
L_a\subset L^2(0,\infty;dt)
\]
is Burnol's extended Sonine space: functions constant on \((0,a)\) whose
cosine transform is again constant on \((0,a)\).  It also places the residual
carrier inside
\[
Y_a\subset L_a.
\]

Step 153 gives the transport and evaluator formula:
\[
T_a:=\MGamma J_a\Uinf^{-1},
\]
and
\[
\Uinf(J_a^*Y^a_{w,k})
=
T_a^*\partial_{\overline w}^{\,k}K_a^\Gamma(\cdot,w).
\]
For \(k=0\), set
\[
\kappa_{a,w}:=T_a^*K_a^\Gamma(\cdot,w).
\]

Step 153 also records the load-bearing projected kernel:
\[
K_a^\Gamma(\cdot,w)
=
P_{\mathcal L_a^\Gamma}
K_a^{\Gamma,\mathrm{amb}}(\cdot,w),
\]
and warns that omitting this projection gives only the ambient Hardy shadow,
not the Burnol/Sonine evaluator.

Step 164 defines
\[
\rho_1=\frac12+14.134725141734693i,\quad
\rho_2=\frac12+21.022039638771554i,\quad
\rho_3=\frac12+25.010857580145688i,
\]
\[
a_0=\frac12,\qquad k=0,
\]
and
\[
e_i=\frac{\eta^{a_0}_{\rho_i,0}}
{\|\eta^{a_0}_{\rho_i,0}\|},
\qquad
\eta^a_{\rho,k}=J_a^*Y^a_{\rho,k}.
\]
The Step 164 matrix entries are
\[
c_{ij}(\ell)=\langle e_i,\Cell e_j\rangle,
\qquad
\Cell=(I-P_\infty)M_{m_\ell}P_\infty,
\]
where
\[
m_\ell(s)=\exp[-\ell(1/2-s)].
\]

\section{T2a. Apply \(\Psym\) on the Mellin Side}

Let \(P=\Psym\), and write \(s=1/2+i\tau\).  Step 173 gives
\[
(P F)(\tau)
=
F(\tau)
-\int_{\mathbb R}
\frac{\sin(\tau-u)}{\pi(\tau-u)}F(u)\,du
-\sum_{n\ge0}\Psi_n(\tau)\langle F,\Psi_n\rangle.
\]
Applying this to the pulled evaluator vector gives
\[
\boxed{
(P\kappa_{a,w})(\tau)
=
\kappa_{a,w}(\tau)
-\int_{\mathbb R}
\frac{\sin(\tau-u)}{\pi(\tau-u)}
\kappa_{a,w}(u)\,du
-\sum_{n\ge0}\Psi_n(\tau)
\langle \kappa_{a,w},\Psi_n\rangle.
}
\]
This formula is explicit once \(\kappa_{a,w}\) is known.

\section{T2b. Transport Back to the Projected Kernel}

The projected Burnol/Sonine kernel entry can be written in the pulled Mellin
model as
\[
\boxed{
K_{\mathrm{proj}}(z,w)
:=
(P_{\mathcal L_a^\Gamma}K_a^\Gamma)(z,w)
=
\langle \kappa_{a,z},P\kappa_{a,w}\rangle.
}
\]
Equivalently, this is the Step 153 transport-back statement:
\[
K_{\mathrm{proj}}(z,w)
=
\left(T_aP\,T_a^{-1}K_a^\Gamma(\cdot,w)\right)(z),
\]
whenever \(T_a^{-1}K_a^\Gamma(\cdot,w)\) is instantiated as
\(\kappa_{a,w}\).

\section{T2c. Specialization to \(\Rho_{\mathrm{fin}}\)}

For the Step 164 finite carrier,
\[
\boxed{
K_{\mathrm{proj}}(\rho_i,\rho_j)
=
\langle \kappa_i,P\kappa_j\rangle,
\qquad
\kappa_i:=\kappa_{a_0,\rho_i}.
}
\]
The norms of finite vectors are
\[
\|\eta^{a_0}_{\rho_i,0}\|=\|\kappa_i\|,
\]
because the Step 153 transport is used through the Hilbert adjoint convention.

\section{T2d. SL164.1 Matrix Entry}

Before normalization,
\[
\langle \eta_i,\Cell\eta_j\rangle
=
\langle \kappa_i,(I-P)M_{m_\ell}P\kappa_j\rangle.
\]
Since \(P\) is an orthogonal projection,
\[
\boxed{
c_{ij}(\ell)
=
\frac{
\langle \kappa_i-P\kappa_i,\,
M_{m_\ell}P\kappa_j\rangle}
{\|\kappa_i\|\,\|\kappa_j\|}.
}
\]
This is the requested operational reduction of SL164.1 to \(K_\infty^{op}\)
plus the Step 153 transport data.

\section{T3. Numerical Attempt}

The requested test pair was
\[
i=j=1,\qquad a_0=\frac12,\qquad
\ell=\log 2=0.6931471805599453.
\]
The computation script recorded the formula
\[
c_{11}(\ell)
=
\frac{\langle \kappa_1-P\kappa_1,M_{m_\ell}P\kappa_1\rangle}
{\|\kappa_1\|^2}.
\]
However, quadrature cannot lawfully begin because the inherited records do not
provide \(\kappa_1(\tau)\).  Step 153 gives only
\[
K_a^\Gamma=P_{\mathcal L_a^\Gamma}K_a^{\Gamma,\mathrm{amb}},
\]
with the projection still load-bearing.  Replacing this by
\(K_a^{\Gamma,\mathrm{amb}}\) would be exactly the public-shadow substitution
forbidden by Step 153.

\[
\boxed{
c_{11}(\log2)\ \text{is not evaluable from inherited records.}
}
\]

\section{Branch B Status}

Step 177 sharpens SL164.1.  The missing record is no longer merely
``projected Sonine kernel unavailable''; it is the explicit Mellin
representative
\[
\kappa_{a,w}=T_a^*K_a^\Gamma(\cdot,w)
\]
or, equivalently, an exact formula for
\[
P_{\mathcal L_a^\Gamma}K_a^{\Gamma,\mathrm{amb}}.
\]
Until that vector is supplied, the finite-carrier matrix diagnostic from Step
164 remains indeterminate.

\section{Nonclaim Boundary}

This step does not prove RH, does not close \(\Xi_{\mathrm{BC}}\), does not
close Branch B, and does not compute a numerical \(c_{ij}(\ell)\).  It does
not weaken retained no-gos and does not use the ambient Hardy shadow as a
substitute for the projected Burnol/Sonine kernel.

\end{document}
